\subsection{半单代数上的导子}

在本小节及随后各小节中，K 是一个交换环，A 是一个 K-代数，B 是 K-代数 \(A \otimes_{K} A^{o}\)，而 \(\varepsilon\) 是由 \(\varepsilon(x \otimes y) = xy\)（对 A 中的 x, y）定义的从 B 到 A 的 K-线性映射。

回顾（III, §4, No.~3, 第 467 页）每个 (A, A)-双模 P 都可看作一个左 B-模，其外部作用律由 \((a \otimes a')p = apa'\)（对 A 中的 \(a, a'\) 与 P 中的 p）给出。反之，每个 B-模都可看作一个 (A, A)-双模。我们赋予 A 它的典范 (A, A)-双模结构及相应的 B-模结构；我们赋予 B 对应于 B-模 \(B_s\) 的 (A, A)-双模结构。于是我们有

\[
a (x \otimes y) a ^ {\prime} = (a \otimes a ^ {\prime}) (x \otimes y) = a x \otimes y a ^ {\prime}
\]

对 A 中的 \(a, a', x, y\) 成立，其中乘积 \(ya'\) 是在代数 A 中计算的。

K-线性映射 \(\varepsilon\) 是 (A, A)-双模的一个同态。

命题 5. —— 下列性质等价：

\begin{enumerate}
\def\labelenumi{(\roman{enumi})}
\item
  B-模 A 是投射的。
\item
  存在 (A,A)-双模 B 的一个元素 \(e\)，满足以下两个条件：\(\varepsilon(e) = 1\) 且 \(ae = ea\)（对每个 \(a \in \mathrm{A}\)）。
\end{enumerate}

映射 \(\varepsilon: \mathrm{B} \to \mathrm{A}\) 是满射的，因为 \(\varepsilon(a \otimes 1) = a\) 对每个 \(a \in \mathrm{A}\) 成立；它是 (A, A)-模的一个同态，因而是 B-线性的。若 B-模 A 是投射的，则存在一个截面 \(s\)（\(\varepsilon\) 的；见 II, §2, No.~2, 第 231 页，命题 4）；它是从 A 到 B 的一个 (A, A)-双模同态。若令 \(e = s(1)\)，则我们有 \(\varepsilon(e) = \varepsilon(s(1)) = 1\) 与 \(ae = s(a) = ea\)（对每个 \(a \in \mathrm{A}\)）。故 (i) 蕴含 (ii)。

反之，设 e 是 B 的一个满足 (ii) 中条件的元素。用下述公式定义从 A 到 B 的一个映射 s：

\[
s (a) = a e = e a.\tag{1}
\]

它是 (A, A)-模的一个同态，且我们有 \(\varepsilon \circ s = 1_{\mathrm{A}}\)；换言之，\(s\) 是满射映射 \(\varepsilon\) 的一个 B-线性截面。因此，B-模 A 同构于子模 \(s(\mathrm{A})\)（\(\mathrm{B}_s\) 的直因子子模；见 II, §1, No.~9, 第 211 页，命题 15），从而是投射的（II, §2, No.~2, 第 231 页，命题 4）。这就证明了 (ii) 蕴含 (i)。

注. —— 1) 设 \(e = \sum_{i=1}^{r} a_i \otimes a_i'\) 是 B 的一个元素。命题 5 的 (ii) 中的条件翻译为公式

\[
\sum_ {i = 1} ^ {r} a _ {i} a _ {i} ^ {\prime} = 1,\tag{2}
\]

\[
\sum_ {i = 1} ^ {r} a a _ {i} \otimes a _ {i} ^ {\prime} = \sum_ {i = 1} ^ {r} a _ {i} \otimes a _ {i} ^ {\prime} a \quad \text { for   every } a \in A.\tag{3}
\]

当它们被满足时，e 是 B 中的一个幂等元。事实上，我们这时有关系式

\[
e ^ {2} = \sum_ {i = 1} ^ {r} a _ {i} e a _ {i} ^ {\prime} = \sum_ {i = 1} ^ {r} e a _ {i} a _ {i} ^ {\prime} = e.
\]

\begin{enumerate}
\def\labelenumi{\arabic{enumi})}
\setcounter{enumi}{1}
\tightlist
\item
  设 K 是一个交换域，A 是一个 K-代数，M 是一个 A-模。加群 \(\operatorname{End}_{\mathrm{K}}(\mathrm{M})\) 被赋予一个由下式定义的 (A, A)-双模结构
\end{enumerate}

\[
a u a ^ {\prime} (x) = a u \left(a ^ {\prime} x\right)
\]

对每个 \(a, a' \in A\)、每个 \(u \in \operatorname{End}_{\mathrm{K}}(\mathrm{M})\) 与每个 \(x \in M\) 成立。我们赋予它相应的 B-模结构。设 \(e = \sum_{i=1}^{r} a_i \otimes a'_i\) 是 B 的一个满足命题 5 中 (ii) 的条件的元素，因而也满足关系式 (2) 与 (3)。若 \(p \in \operatorname{End}_{\mathrm{K}}(\mathrm{M})\) 是一个投影，其像 N 是 M 的一个 A-子模，则 ep 是具有同一像的一个 A-线性投影。

事实上，\(ep\) 的像包含在 N 中。若 \(x\) 属于 N，则 \(a_i' x\) 亦然，于是我们有 \(p(a_i' x) = a_i' x\) 且

\[
e p (x) = \sum_ {i = 1} ^ {r} a _ {i} a _ {i} ^ {\prime} x = x
\]

这由公式 (2) 得出。由公式 (3) 我们推出，\(aep(x) = ep(ax)\) 对每个 \(a \in \mathrm{A}\) 与每个 \(x \in \mathrm{M}\) 成立，这就证明了 \(ep\) 是 A-线性的。

定理 2. —— 设 K 是一个交换域，A 是一个 K-代数。下列性质等价：

\begin{enumerate}
\def\labelenumi{(\roman{enumi})}
\item
  K-代数 A 是绝对半单的。
\item
  K-代数 \(\mathrm{B} = \mathrm{A}\otimes_{\mathrm{K}}\mathrm{A}^{\circ}\) 是半单的。
\item
  B-模 A 是投射的。
\item
  存在 (A,A)-双模 B 的一个元素 \(e\)，满足 \(\varepsilon(e) = 1\) 且 \(ae = ea\)（对每个 \(a \in \mathrm{A}\)）。
\end{enumerate}

假设代数 A 是绝对半单的，因而是半单的。则代数 \(A^{o}\) 是半单的（VIII, 第 137 页，命题 2），且由 VIII, 第 234 页的推论 1 得出，\(B = A \otimes_{K} A^{o}\) 是一个半单 K-代数。因此，(i) 蕴含 (ii)。

由于半单环上的每个模都是投射的（VIII, 第 138 页，命题 4），(ii) 蕴含 (iii)。

(iii) 与 (iv) 的等价性由命题 5 得出。为完成证明，我们来证明 (iv) 蕴含 (i)。设 \(e = \sum_{i=1}^{r} a_i \otimes a_i'\) 是 B 的一个满足命题 5 中 (ii) 的条件的元素。设 L 是域 K 的一个扩张；我们必须证明环 \(\mathrm{A}_{(\mathrm{L})}\) 是半单的，或者等价地，每个 \(\mathrm{A}_{(\mathrm{L})}\) -模都是半单的（VIII, 第 138 页，命题 4）。设 M 是一个 \(\mathrm{A}_{(\mathrm{L})}\) -模，N 是 M 的一个子模；我们把 M 看作左 (A,L)-双模，把 N 看作一个子双模（III, §4, No.~3, 第 466 页）。由于 L 是一个域，存在 M 中以 N 为像的 L-线性投影 \(u\)。由于位似 \(a_{\mathrm{M}}\)（与 A 的元素 \(a\) 相伴的）都是 L-线性的，存在从 \(\mathrm{A} \otimes_{\mathrm{K}} \mathrm{A}^{\circ}\) 到 \(\operatorname{End}_{\mathrm{L}}(\mathrm{M})\) 的唯一的加群同态，它把元素 \(a \otimes a'\) 映为 L-线性映射 \(x \mapsto au(a'x)\)。我们把此同态下 \(e\) 的像记作 \(v\)；由注 2 得出，\(v\) 是一个以 N 为像的 \(\mathrm{A}_{(\mathrm{L})}\) -线性投影。\(v\) 的核是 M 的一个 \(\mathrm{A}_{(\mathrm{L})}\) -子模，它与 N 互补。由 VIII, 第 56 页的推论 2，\(\mathrm{A}_{(\mathrm{L})}\) -模 M 是半单的。

注 3. —— 我们知道（VIII, 第 234 页，推论 1），交换域上的两个绝对半单代数的张量积是绝对半单的。因此，若代数 A 是绝对半单的，则代数 \(B = A \otimes_{K} A^{\circ}\) 亦然。

